\subsubsection{Orthogonale Projektion}
\label{subsubsec:vector_algebra_orthogonal_projection}

Mit dem Skalarprodukt lässt sich die orthogonale Projektion
\EnglishTerm{orthogonal projection} eines Vektors auf einen anderen
berechnen. Die unterschiedliche Wirkung der Komponenten wird an einer
gelenkig gelagerten Wippe unmittelbar sichtbar: Die parallele Komponente
\EnglishTerm{parallel component} belastet das Brett in Längsrichtung, die
senkrechte Komponente \EnglishTerm{orthogonal component} erzeugt ein
Drehmoment um das Gelenk.

\begin{center}
  \CachedTikzFigure{content/resources/chapters/vector_algebra/tikz/orthogonal_projection_seesaw.tex}
\end{center}

Um diese Wirkungen getrennt untersuchen zu können, zerlegen wir einen Vektor
$\vec b$ relativ zu einer Richtung $\vec a\ne\vec 0$ in zwei Komponenten. Für diese Zerlegung gilt:
\par\medskip
\noindent
\begin{minipage}[c]{0.45\linewidth}
  \begin{itemize}[leftmargin=1.4em,itemsep=0.8em]
    \item Parallel zur Richtung:\quad
      $\vec b_{\parallel}\parallel\vec a$
    \item Senkrecht zur Richtung:\quad
      $\vec b_{\perp}\perp\vec a$
    \item Vollständige Zerlegung:\quad
      $\vec b=\vec b_{\parallel}+\vec b_{\perp}$
  \end{itemize}
\end{minipage}\hfill
\begin{minipage}[c]{0.52\linewidth}
  \centering
  \CachedTikzFigure{content/resources/chapters/vector_algebra/tikz/orthogonal_projection.tex}
\end{minipage}
\par\medskip

\paragraph{Parallele Komponente}
Zunächst sei der Zwischenwinkel $\varphi<90^\circ$. Dann zeigen
$\vec b_{\parallel}$ und $\vec a$ in dieselbe Richtung.
Für die \emph{Länge} der parallelen Komponente nutzen wir zwei Regeln des
Skalarprodukts: Bei gleichgerichteten Vektoren ist es das Produkt der
Beträge, bei senkrechten Vektoren ist es null.
\begin{ClozeArea}[34]{8}
  \[
    \begin{aligned}
      \vec a\cdot\vec b
        &=\vec a\cdot\bigl(\vec b_{\parallel}+\vec b_{\perp}\bigr)
         =\underbrace{\vec a\cdot\vec b_{\parallel}}_{
            \lvert\vec a\rvert\,\lvert\vec b_{\parallel}\rvert}
          +\underbrace{\vec a\cdot\vec b_{\perp}}_{0}, \\[1em]
      \Longrightarrow\quad \lvert\vec b_{\parallel}\rvert
        &=\vec b\cdot\frac{\vec a}{\lvert\vec a\rvert}
         =\vec b\cdot\vec e_a.
    \end{aligned}
  \]
\end{ClozeArea}
Ihre \emph{Richtung} gibt der Einheitsvektor $\vec e_a$ vor. Länge mal
Einheitsvektor ergibt damit die parallele Komponente.
\begin{ClozeArea}[34]{4}
  \[
    \vec b_{\parallel}
      =\lvert\vec b_{\parallel}\rvert\,\vec e_a
      =\bigl(\vec b\cdot\vec e_a\bigr)\,\vec e_a.
  \]
\end{ClozeArea}

Bei $\varphi>90^\circ$ können wir dieselbe Projektionsgerade durch
$-\vec a$ beschreiben und erhalten einen Zwischenwinkel kleiner als
$90^\circ$. In der Formel wechseln dadurch $\vec b\cdot\vec e_a$ und
$\vec e_a$ beide das Vorzeichen, ihr Produkt bleibt gleich. Bei
$\varphi=90^\circ$ oder $\vec b=\vec 0$ liefert die Formel
$\vec b_{\parallel}=\vec 0$.

Die senkrechte Komponente erhalten wir durch Umstellen der ursprünglichen
Zerlegung und anschließendes Einsetzen von $\vec b_{\parallel}$:
\begin{ClozeArea}[34]{4}
  \[
    \vec b=\vec b_{\parallel}+\vec b_{\perp}
      \quad\Longleftrightarrow\quad
    \vec b_{\perp}
      =\vec b-\vec b_{\parallel}
  \]
\end{ClozeArea}

\begin{TheoremBox}{Orthogonale Projektion}{vector_algebra_projection_components}
  Projiziert man $\vec b\in\R^n$ orthogonal auf die durch
  $\vec a\in\R^n\setminus\{\vec 0\}$ vorgegebene Richtung, so erhält man die
  Zerlegung
  $\textcolor{red!75!black}{\vec b}=\textcolor{black}{\vec b_{\parallel}}+\textcolor{olive!80!black}{\vec b_{\perp}}$
  für die
  \par\medskip
  \noindent
  \begin{minipage}[c]{0.54\linewidth}
    \begin{itemize}[leftmargin=1.5em,itemsep=0.8em]
      \item \emph{parallele Komponente}
        $\vec b_{\parallel}$ mit $\vec b_{\parallel}\parallel\vec a$:
        \[
          \textcolor{black}{\vec b_{\parallel}}
          =\bigl(\textcolor{red!75!black}{\vec b}\cdot
            \textcolor{THMBlue}{\vec e_a}\bigr)
            \textcolor{THMBlue}{\vec e_a}.
        \]

      \item \emph{senkrechte Komponente}
        $\vec b_{\perp}$ mit $\vec b_{\perp}\perp\vec a$:
        \[
          \textcolor{olive!80!black}{\vec b_{\perp}}
            =\textcolor{red!75!black}{\vec b}
            -\textcolor{black}{\vec b_{\parallel}}.
        \]
    \end{itemize}
  \end{minipage}\hfill
  \begin{minipage}[c]{0.43\linewidth}
    \centering
    \CachedTikzFigure{content/resources/chapters/vector_algebra/tikz/orthogonal_projection.tex}
  \end{minipage}
\end{TheoremBox}

\begin{ExampleBox}{Einen Vektor orthogonal zerlegen}{vector_algebra_projection_example}
  \noindent
  \begin{minipage}[c]{0.60\linewidth}
    Gegeben seien
    \[
      \vec a=\begin{pmatrix}4\\2\end{pmatrix},
      \qquad
      \vec b=\begin{pmatrix}1\\3\end{pmatrix}.
    \]
    Zunächst bestimmen wir den Einheitsvektor der Projektionsrichtung:
    \[
      \vec e_a
        =\frac{\vec a}{\lvert\vec a\rvert}
        =\frac{1}{\sqrt{20}}\begin{pmatrix}4\\2\end{pmatrix}
        =\frac{1}{\sqrt{5}}\begin{pmatrix}2\\1\end{pmatrix}.
    \]
    Anschließend berechnen wir $\vec b_{\parallel}$ und $\vec b_{\perp}$:
    \[
      \begin{aligned}
        \vec b_{\parallel}
          &=(\vec b\cdot\vec e_a)\vec e_a
            =\frac{5}{\sqrt{5}}\,
            \frac{1}{\sqrt{5}}\begin{pmatrix}2\\1\end{pmatrix}
            =\Cloze{2.1cm}{\begin{pmatrix}2\\1\end{pmatrix}}, \\[0.8em]
        \vec b_{\perp}
          &=\vec b-\vec b_{\parallel}
            =\begin{pmatrix}1\\3\end{pmatrix}
            -\begin{pmatrix}2\\1\end{pmatrix}
            =\Cloze{2.1cm}{\begin{pmatrix}-1\\2\end{pmatrix}}.
      \end{aligned}
    \]
    Zur Kontrolle eignen sich zwei kurze Rechnungen:
    \[
      \begin{array}{@{}r@{\qquad}l@{}}
        \text{Orthogonalität:}
          &\displaystyle \vec a\cdot\vec b_{\perp}
            =\begin{pmatrix}4\\2\end{pmatrix}\cdot
            \begin{pmatrix}-1\\2\end{pmatrix}
            =\Cloze{3cm}{4\cdot(-1)+2\cdot2=0}, \\[0.8em]
        \text{Vollständigkeit:}
          &\displaystyle \vec b_{\parallel}+\vec b_{\perp}
            =\begin{pmatrix}2\\1\end{pmatrix}
            +\begin{pmatrix}-1\\2\end{pmatrix}
            =\Cloze{3cm}{\begin{pmatrix}1\\3\end{pmatrix}=\vec b}.
      \end{array}
    \]
  \end{minipage}\hfill
  \begin{minipage}[c]{0.38\linewidth}
    \centering
    \CachedTikzFigure{content/resources/chapters/vector_algebra/tikz/orthogonal_projection_example.tex}
  \end{minipage}
\end{ExampleBox}
